The first image of a black hole was constructed recently by the international
team of the Event Horizon Telescope (EHT). The imaged area surrounds the
supermassive black hole in the centre of the galaxy M87. The observations
resulting in the final image were carried out at a wavelength
$\lambda = 1.3$\,mm where the interstellar extinction is not prohibitively
large.

\begin{description}
\item[a)] How large an instrument would be needed to resolve the shadow (in
  effect the photon capture radius, which is three times the size of event
  horizon) of a supermassive black hole in the centre of a galaxy? Express the
  result as a function of the distance $d$ and the mass $M$ of the black hole.
  \hfill[6\,p]
\item[b)] Give the size of the instrument in units of Earth radius for
  \begin{description}
  \item[i.] the supermassive black hole in the centre of M87
    ($d_{\mathrm{BH\text{-}M87}} = 5.5 \times 10^7$\,ly,
    $M_{\mathrm{BH\text{-}M87}} = 6.5 \times 10^9\,\mathrm{M}_\odot$),
    \hfill[1\,p]
  \item[ii.] and Sgr~A*, the supermassive black hole of our own galaxy, Milky
    Way ($d_{\mathrm{Sgr\,A^*}} = 8.3$\,kpc,
    $M_{\mathrm{Sgr\,A^*}} = 3.6 \times 10^6\,\mathrm{M}_\odot$). \hfill[1\,p]
  \end{description}
\item[c)] What type of technology is needed for the development of such an
  instrument? Mark the letter of your answer with $\times$ on the answer
  sheet. \hfill[2\,p]
  \begin{description}
  \item[(A)] Gravitational lensing by dark matter
  \item[(B)] Interferometry with an array of radio telescopes
  \item[(C)] Photon deceleration in a dense environment
  \item[(D)] Reducing the effect of incoming wavefront distortions
  \item[(E)] Neutrino focusing with strong electromagnetic fields
  \end{description}
\end{description}
